\documentclass[10pt]{extarticle}
\usepackage{fontspec}
\setmainfont{Leelawadee UI}
\XeTeXlinebreaklocale "th_TH"
\XeTeXlinebreakskip = 0pt plus 1pt minus 0.1pt
\usepackage[a4paper,top=1.3cm,bottom=1.5cm,left=1.4cm,right=1.4cm]{geometry}
\usepackage{amsmath,amssymb}
\usepackage{xcolor}
\usepackage{colortbl}
\usepackage{booktabs}
\usepackage{tikz}
\usepackage{pgfplots}
\pgfplotsset{compat=1.18}
\usepackage{enumitem}
\setlist{nosep,leftmargin=1.3em,itemsep=1pt,topsep=2pt}
\usepackage{titlesec}
\usepackage{tcolorbox}
\tcbuselibrary{skins,breakable}
\usepackage{array}
\usepackage{longtable}
\usepackage{multirow}

\definecolor{navy}{RGB}{20,40,90}
\definecolor{gold}{RGB}{190,140,20}
\definecolor{lightbg}{RGB}{240,244,250}
\definecolor{solbg}{RGB}{247,247,240}
\definecolor{ideabg}{RGB}{236,246,238}
\definecolor{ideafr}{RGB}{60,120,70}
\definecolor{compbg}{RGB}{236,247,247}
\definecolor{compfr}{RGB}{20,120,120}
\definecolor{qbg}{RGB}{245,240,250}
\definecolor{qfr}{RGB}{90,60,140}

\titleformat{\section}{\normalfont\bfseries\Large\color{navy}}{\thesection}{0.5em}{}[{\color{gold}\titlerule[1.2pt]}]
\titlespacing*{\section}{0pt}{10pt}{5pt}
\titleformat{\subsection}{\normalfont\bfseries\large\color{gold!70!black}}{\thesubsection}{0.4em}{}
\titlespacing*{\subsection}{0pt}{7pt}{3pt}
\titleformat{\subsubsection}{\normalfont\bfseries\color{navy}}{\thesubsubsection}{0.4em}{}
\titlespacing*{\subsubsection}{0pt}{5pt}{2pt}

\newtcolorbox{formulabox}{
  colback=lightbg, colframe=navy!60, boxrule=0.4pt, breakable,
  arc=1pt, left=4pt, right=4pt, top=3pt, bottom=3pt,
}
\newtcolorbox{solbox}{
  colback=solbg, colframe=gold!70!black, boxrule=0.4pt, breakable,
  arc=1pt, left=4pt, right=4pt, top=3pt, bottom=3pt,
}
\newtcolorbox{ideabox}{
  colback=ideabg, colframe=ideafr, boxrule=0.4pt, breakable,
  arc=1pt, left=4pt, right=4pt, top=3pt, bottom=3pt,
}
\newtcolorbox{outbox}{
  colback=compbg, colframe=compfr, boxrule=0.4pt, breakable,
  arc=1pt, left=4pt, right=4pt, top=3pt, bottom=3pt,
}
\newtcolorbox{qbox}{
  colback=qbg, colframe=qfr, boxrule=0.5pt, breakable,
  arc=1pt, left=4pt, right=4pt, top=3pt, bottom=3pt,
  fonttitle=\bfseries\color{white}, colbacktitle=qfr, title={โจทย์}
}
\newtcolorbox{stepbox}[1]{
  colback=white, colframe=navy, boxrule=0.6pt, breakable,
  arc=1pt, left=5pt, right=5pt, top=3pt, bottom=4pt,
  fonttitle=\bfseries\color{white}, colbacktitle=navy, title={#1}
}

\newcommand{\dist}[1]{\subsection*{\textcolor{navy}{\textbullet\ #1}}}
\newcommand{\case}[1]{\textbf{\textcolor{gold!70!black}{#1}}}
\newcommand{\concept}[1]{\par\vspace{2pt}\noindent{\footnotesize\textbf{\textcolor{ideafr}{\ding{228}\ แนวคิด:}} #1}\vspace{2pt}\par}
\newcommand{\ding}[1]{$\blacktriangleright$}
\newcommand{\spssout}[1]{\vspace{2pt}\noindent{\small\textbf{\textcolor{compfr}{Output -- #1}}}\par\vspace{1pt}}

\setlength{\parindent}{0pt}
\setlength{\parskip}{1pt}
\renewcommand{\arraystretch}{1.2}

\pgfmathdeclarefunction{gauss}{3}{%
  \pgfmathparse{1/(#3*sqrt(2*pi))*exp(-((#1-#2)^2)/(2*#3^2))}%
}

\begin{document}
\begin{center}
{\LARGE\bfseries\color{navy} Chapter 4: การประมาณค่า (Estimation)}\\[3pt]
{\normalsize สรุปเนื้อหาพร้อมแนวคิดและวิธีทำละเอียด — อ.ดร.ธนวรรณ ประฮาดไชย \quad Department of Statistics, Faculty of Science, KKU}
\end{center}
\vspace{3pt}
{\color{gold}\hrule height 1.6pt}
\vspace{6pt}

\section{แนวคิดพื้นฐานของการประมาณค่า}

บทนี้ตอบคำถามว่า \textbf{เมื่อเราไม่ทราบค่าพารามิเตอร์ที่แท้จริงของประชากร แล้วมีแต่ข้อมูลจากตัวอย่างสุ่ม เราจะ ``เดา'' ค่าพารามิเตอร์นั้นอย่างเป็นระบบและน่าเชื่อถือได้อย่างไร} หลักการคือ ใช้ข้อมูลจาก\textbf{ตัวอย่าง} (สุ่มมาจากประชากร) มาประมาณค่า\textbf{พารามิเตอร์}ของประชากรที่ไม่ทราบค่าจริง แล้วรายงานว่าค่าที่ประมาณได้นั้นมีความ ``น่าเชื่อถือ'' เพียงใด

\begin{center}
\begin{tabular}{@{}ll@{}}
\rowcolor{lightbg}\textbf{ประชากร (Population)} & \textbf{ตัวอย่าง (Sample)}\\
พารามิเตอร์ (Parameter) & ตัวสถิติ (Statistic)\\
ค่าเฉลี่ย: $\mu$ & ค่าเฉลี่ย: $\bar{X}$\\
ค่าสัดส่วน: $p$ & ค่าสัดส่วน: $\hat{p}$\\
\end{tabular}
\end{center}

\dist{คำศัพท์ที่ต้องแยกให้ออก}
\textbf{พารามิเตอร์ ($\theta$):} ค่าคงที่ที่แสดงลักษณะบางอย่างของประชากรซึ่ง\textbf{ไม่ทราบค่าที่แท้จริง} เช่น ค่าเฉลี่ยประชากร ค่าสัดส่วนประชากร ความแปรปรวนประชากร

\textbf{ตัวสถิติ ($\hat\theta$):} ค่าที่คำนวณได้จากตัวอย่างสุ่ม (หรือฟังก์ชันของตัวอย่างสุ่ม) ซึ่งมีค่า\textbf{แตกต่างกันไปตามตัวอย่างที่สุ่มได้แต่ละครั้ง} เนื่องจากมีตัวอย่างที่สุ่มได้จำนวนมากมาย จึงถือว่าตัวสถิติเป็น\textbf{ตัวแปรสุ่ม}เช่นเดียวกัน

\textbf{ตัวประมาณ (Estimator):} ตัวสถิติที่นำไปใช้ประมาณค่าพารามิเตอร์ (คือ ``สูตร'' หรือ ``กติกา'' ที่จะใช้คำนวณ)

\textbf{ค่าประมาณ (Estimate):} ค่าตัวเลขจริงของตัวสถิตินั้นที่คำนวณได้จากตัวอย่างที่สุ่มมา (คือ ``ผลลัพธ์'' ที่ได้จากการแทนค่าตามสูตร)

\begin{qbox}
นักวิจัยต้องการศึกษา\textbf{ค่าใช้จ่ายเฉลี่ยต่อเดือน}ของนักศึกษามหาวิทยาลัยขอนแก่น ซึ่งมีทั้งสิ้น $N=40{,}000$ คน เนื่องจากไม่สามารถเก็บข้อมูลนักศึกษาทั้ง 40,000 คนได้ (เสียเวลาและค่าใช้จ่ายมาก) จึงสุ่มตัวอย่างนักศึกษามา $n=500$ คน ได้ค่าใช้จ่ายเฉลี่ยเท่ากับ 5{,}000 บาท
\end{qbox}
\begin{ideabox}
\textbf{แนวคิด:} $\mu$ คือ\textbf{พารามิเตอร์} แทนค่าใช้จ่ายเฉลี่ยต่อเดือนที่แท้จริงของนักศึกษาทั้ง 40,000 คน (ไม่ทราบค่า); $\bar{X}$ คือ\textbf{ตัวประมาณ} (สูตร $\bar X=\sum X_i/n$ ที่จะใช้คำนวณจากตัวอย่าง 500 คน); และ $\bar{X}=5{,}000$ บาท คือ\textbf{ค่าประมาณ} ของ $\mu$ (ตัวเลขที่คำนวณได้จริง)
\end{ideabox}

\dist{คุณสมบัติของตัวประมาณที่ดี}
มีเกณฑ์ 3 ข้อที่ใช้ตัดสินว่าตัวประมาณหนึ่ง ``ดี'' เพียงใด:

\textbf{1) ไม่เอนเอียง (Unbiased estimator):} ถ้า $E(\hat\theta)=\theta$ เรียก $\hat\theta$ ว่าเป็นตัวประมาณที่ไม่เอนเอียงของ $\theta$ กล่าวคือ ถ้าสุ่มตัวอย่างซ้ำ ๆ กันหลายครั้งแล้วคำนวณ $\hat\theta$ ทุกครั้ง ค่าเฉลี่ยของ $\hat\theta$ ทั้งหมดจะเท่ากับค่าพารามิเตอร์ที่แท้จริงพอดี (ไม่สูงหรือต่ำกว่าอย่างเป็นระบบ)

\textbf{2) ไม่เอนเอียงและมีความแปรปรวนต่ำสุด (มีประสิทธิภาพ):} สมมติ $\bar X$ และ $Me$ (มัธยฐานตัวอย่าง) ต่างก็เป็นตัวประมาณที่ไม่เอนเอียงของ $\mu$ ทั้งคู่ ถ้าพิสูจน์ได้ว่า $Var(\bar X) < Var(Me)$ แสดงว่า $\bar X$ มีประสิทธิภาพมากกว่า $Me$ (ค่าประมาณที่ได้จาก $\bar X$ จะกระจายรอบ $\mu$ แคบกว่า เชื่อถือได้มากกว่า) จึงเลือกใช้ $\bar X$

\textbf{3) ความแนบนัย/ความคงเส้นคงวา (Consistency):} เมื่อขนาดตัวอย่าง $n$ เพิ่มขึ้นเข้าใกล้ขนาดประชากร $N$ (หรือเพิ่มขึ้นไม่จำกัด) ค่าประมาณ $\hat\theta$ จะเข้าใกล้ค่าพารามิเตอร์ที่แท้จริงมากขึ้นเรื่อย ๆ เช่น $\bar X \to \mu$ เมื่อ $n\to N$

\dist{2 วิธีการประมาณค่าพารามิเตอร์}
\case{1) การประมาณแบบจุด (Point estimation):} ใช้ข้อมูลตัวอย่างคำนวณ\textbf{ค่าตัวเลขเดียว}เป็นค่าประมาณของพารามิเตอร์ เช่น $\bar{X}$ (ประมาณ $\mu$), $\hat{p}=x/n$ (ประมาณ $p$), $S^2$ (ประมาณ $\sigma^2$) -- \textit{ข้อจำกัด:} ค่าประมาณแบบจุดมีโอกาสคลาดเคลื่อนไปจากพารามิเตอร์จริงได้เสมอ (จากวิธีสุ่มตัวอย่าง ขนาดตัวอย่าง และการกระจายของข้อมูล) แต่ตัวเลขเพียงค่าเดียวไม่ได้บอกว่า ``คลาดเคลื่อนได้มากแค่ไหน''

\case{2) การประมาณแบบช่วง (Interval estimation):} เพื่อแก้ข้อจำกัดข้างต้น จึงหา\textbf{ช่วง} $(L,U)$ ที่มีความน่าจะเป็นระดับหนึ่งว่าจะครอบคลุมพารามิเตอร์ที่แท้จริง แทนที่จะรายงานตัวเลขเดียว

\begin{formulabox}
$$P(L<\theta<U)=1-\alpha$$
เรียกช่วง $(L,U)$ ว่า \textbf{ช่วงความเชื่อมั่น} (Confidence Interval) ที่ระดับ $(1-\alpha)\times100\%$ ของ $\theta$ โดย
\begin{itemize}
\item $L$ = ขอบเขตล่างของช่วง (Lower bound), $U$ = ขอบเขตบนของช่วง (Upper bound)
\item $\alpha$ = ระดับนัยสำคัญ (ความน่าจะเป็นที่ช่วงจะ \textit{ไม่} ครอบคลุมพารามิเตอร์จริง)
\item $1-\alpha$ = \textbf{ระดับความเชื่อมั่น} (Confidence level) นิยมกำหนดที่ 0.90, 0.95 หรือ 0.99 (คือ 90\%, 95\%, 99\%)
\end{itemize}
\textbf{ความหมายที่ถูกต้อง:} ถ้าสุ่มตัวอย่างซ้ำ ๆ กันหลายครั้ง แล้วสร้างช่วงความเชื่อมั่นแบบเดียวกันทุกครั้ง จะมีสัดส่วน $(1-\alpha)\times100\%$ ของช่วงทั้งหมดที่สร้างขึ้นครอบคลุมค่าพารามิเตอร์จริง (ไม่ใช่ว่า ``มีโอกาส $1-\alpha$ ที่ $\theta$ จะอยู่ในช่วงที่คำนวณได้ครั้งนี้'' เพราะ $\theta$ เป็นค่าคงที่ ไม่ใช่ตัวแปรสุ่ม)
\end{formulabox}

\begin{center}
\begin{tikzpicture}
\begin{axis}[width=4.6cm,height=2.4cm,axis lines=left,
  xlabel={},ylabel={},ymin=0,ymax=0.45,
  xtick={-2,0,2},xticklabels={$L$,$\theta$,$U$},tick label style={font=\tiny},
  label style={font=\tiny},axis line style={-latex,thin}]
\addplot[domain=-3.2:3.2,samples=70,navy,thick] {gauss(x,0,1)};
\addplot[domain=-2:2,samples=40,fill=gold!45,draw=none] {gauss(x,0,1)} \closedcycle;
\end{axis}
\end{tikzpicture}
\end{center}

\textbf{ความกว้างของช่วงความเชื่อมั่นขึ้นอยู่กับ 3 ปัจจัย:} (1) ความแปรปรวนของข้อมูล ($\sigma^2$ หรือ $S^2$) -- ข้อมูลกระจายมากยิ่งทำให้ช่วงกว้าง (2) ขนาดตัวอย่าง $n$ -- $n$ ยิ่งมากยิ่งทำให้ช่วงแคบลง (แม่นยำขึ้น) (3) ระดับความเชื่อมั่นที่ต้องการ -- ยิ่งต้องการความเชื่อมั่นสูง (เช่น 99\%) ช่วงยิ่งต้องกว้างขึ้นเพื่อให้ ``ครอบคลุม'' ได้บ่อยตามที่ต้องการ

\textbf{ตารางค่าวิกฤต $z_\alpha$ ที่ใช้บ่อย} (นิยาม $P(Z\ge z_\alpha)=\alpha$; ในการประมาณช่วงแบบสองด้าน ให้นำ $\alpha/2$ ไปเปิดตารางแทน $\alpha$):
\begin{center}
\small
\begin{tabular}{r|cccccc}
\rowcolor{lightbg}$\alpha$ & .10 & .05 & .025 & .01 & .005 & .001\\\hline
$z_\alpha$ & 1.2816 & 1.6449 & 1.9600 & 2.3263 & 2.5758 & 3.0902\\
\end{tabular}
\end{center}
ตัวอย่างการอ่านตาราง: ความเชื่อมั่น 95\% $\Rightarrow \alpha=0.05 \Rightarrow \alpha/2=0.025 \Rightarrow z_{0.025}=1.96$; ความเชื่อมั่น 90\% $\Rightarrow \alpha/2=0.05 \Rightarrow z_{0.05}=1.645$; ความเชื่อมั่น 99\% $\Rightarrow \alpha/2=0.005 \Rightarrow z_{0.005}=2.576$

\section{ประมาณค่าเฉลี่ยของ 1 ประชากร ($\mu$)}

\dist{การประมาณแบบจุด}
\begin{formulabox}
$$\bar{X}=\frac{\sum_{i=1}^n X_i}{n} \quad \text{เป็นตัวประมาณแบบจุดของ } \mu$$
\end{formulabox}

\begin{qbox}
สุ่มนักศึกษา 8 คน วัดเวลาที่ใช้ทำข้อสอบ (นาที): 72, 68, 75, 70, 74, 69, 71, 73 จงหาค่าประมาณแบบจุดของ $\mu$
\end{qbox}
\concept{โจทย์ขอเพียง ``ค่าประมาณแบบจุด'' คือค่าเดียว ไม่ได้ถามช่วงความเชื่อมั่น จึงใช้สูตร $\bar X$ ตรง ๆ โดยไม่ต้องพิจารณาว่าทราบ $\sigma$ หรือไม่ และไม่ต้องมีระดับความเชื่อมั่นเข้ามาเกี่ยวข้อง}
\begin{solbox}
$$\bar{X}=\frac{72+68+75+70+74+69+71+73}{8}=\frac{572}{8}=71.5 \text{ นาที}$$
\end{solbox}

\dist{การประมาณแบบช่วง — 3 กรณี}
\begin{formulabox}
\case{กรณี 1 (ทราบ $\sigma^2$ ของประชากร; $n$ จะมากหรือน้อยก็ได้):}
$$\bar{X}\pm z_{\frac{\alpha}{2}}\dfrac{\sigma}{\sqrt{n}}$$
\case{กรณี 2 (ไม่ทราบ $\sigma^2$, $n<30$; ต้องมีข้อสมมติว่าประชากรแจกแจงปกติ):}
$$\bar{X}\pm t_{\frac{\alpha}{2}}\dfrac{S}{\sqrt{n}},\quad v=n-1$$
\case{กรณี 3 (ไม่ทราบ $\sigma^2$, $n\ge30$; ใช้ $S$ แทน $\sigma$ ได้เพราะทฤษฎีลิมิตกลาง):}
$$\bar{X}\pm z_{\frac{\alpha}{2}}\dfrac{S}{\sqrt{n}}$$
โดย $\displaystyle S=\sqrt{\dfrac{\sum(X_i-\bar{X})^2}{n-1}}$ คือส่วนเบี่ยงเบนมาตรฐานของ\textbf{ตัวอย่าง} (ตัวประมาณของ $\sigma$)
\end{formulabox}

\textbf{ตัวแปรสำคัญ}: $\bar X$ = ค่าเฉลี่ยตัวอย่าง (จุดศูนย์กลางของช่วง), $\dfrac{\sigma}{\sqrt n}$ หรือ $\dfrac{S}{\sqrt n}$ = ค่าคลาดเคลื่อนมาตรฐาน (Standard Error) ของ $\bar X$ คือ ส่วนเบี่ยงเบนมาตรฐานของการแจกแจงของ $\bar X$ เอง, $z_{\alpha/2}$ หรือ $t_{\alpha/2}$ = ค่าวิกฤตที่กำหนดความกว้างของช่วงตามระดับความเชื่อมั่นที่ต้องการ, $v=n-1$ = องศาอิสระ (degrees of freedom) ของการแจกแจง $t$

\textbf{ที่มาของสูตร (กรณี 1):} เนื่องจากสุ่มตัวอย่างขนาด $n$ จากประชากรที่แจกแจงปกติ ค่าเฉลี่ย $\mu$ (ไม่ทราบค่า) ความแปรปรวน $\sigma^2$ (ทราบค่า) จะได้ว่า
$$\bar X \sim N\!\left(\mu,\frac{\sigma^2}{n}\right) \quad\Rightarrow\quad Z=\frac{\bar{X}-\mu}{\sigma/\sqrt n}\sim N(0,1)$$
จากสมบัติของโค้งปกติมาตรฐาน พื้นที่ตรงกลาง $1-\alpha$ อยู่ระหว่าง $-z_{\alpha/2}$ กับ $z_{\alpha/2}$:
$$P\!\left(-z_{\frac{\alpha}{2}}<Z<z_{\frac{\alpha}{2}}\right)=1-\alpha$$
แทน $Z$ ด้วยสูตร แล้ว\textbf{จัดรูปสมการทีละขั้น}เพื่อแยก $\mu$ ไว้ตรงกลาง:
$$P\!\left(-z_{\frac{\alpha}{2}}<\frac{\bar{X}-\mu}{\sigma/\sqrt n}<z_{\frac{\alpha}{2}}\right)=1-\alpha$$
คูณทุกส่วนด้วย $\sigma/\sqrt n$:
$$P\!\left(-z_{\frac{\alpha}{2}}\frac{\sigma}{\sqrt n}<\bar{X}-\mu<z_{\frac{\alpha}{2}}\frac{\sigma}{\sqrt n}\right)=1-\alpha$$
ลบ $\bar X$ ออกจากทุกส่วน แล้วคูณด้วย $-1$ (สลับเครื่องหมายอสมการ):
$$P\!\left(\bar{X}-z_{\frac{\alpha}{2}}\frac{\sigma}{\sqrt n}<\mu<\bar{X}+z_{\frac{\alpha}{2}}\frac{\sigma}{\sqrt n}\right)=1-\alpha$$
จึงได้ช่วงความเชื่อมั่น $(1-\alpha)\times100\%$ ของ $\mu$ ตามสูตรกรณี 1 (กรณี 2 และ 3 ได้จากแนวคิดเดียวกัน เพียงเปลี่ยน $\sigma$ เป็น $S$ และเปลี่ยนการแจกแจงของ $Z$ เป็น $T$ เมื่อไม่ทราบ $\sigma$ และตัวอย่างมีขนาดเล็ก)

\begin{qbox}
น้ำหนักสุทธิของน้ำตาลที่บรรจุถุงโดยเครื่องจักรเครื่องหนึ่ง มีการแจกแจงปกติที่มีค่าเบี่ยงเบนมาตรฐาน $\sigma=1.2$ ออนซ์ สุ่มน้ำตาลมา 25 ถุง พบว่ามีน้ำหนักสุทธิเฉลี่ย $\bar X=19.8$ ออนซ์ จงหาความเชื่อมั่น 95\% ของน้ำหนักสุทธิเฉลี่ยของน้ำตาลที่บรรจุถุงโดยเครื่องจักรเครื่องนี้ทั้งหมด
\end{qbox}

\begin{stepbox}{ขั้นที่ 1: ทำความเข้าใจโจทย์}
โจทย์บอกว่าประชากร (น้ำหนักน้ำตาลทุกถุงที่บรรจุโดยเครื่องนี้) แจกแจงแบบปกติ และ\textbf{ทราบค่าเบี่ยงเบนมาตรฐานของประชากร} $\sigma=1.2$ ออนซ์ มาแล้ว ส่วนตัวอย่างสุ่มมา $n=25$ ถุง คำนวณได้ $\bar X=19.8$ ออนซ์ สิ่งที่ต้องการคือ\textbf{ช่วงความเชื่อมั่นของค่าเฉลี่ยประชากร} $\mu$ (น้ำหนักสุทธิเฉลี่ยที่แท้จริงของน้ำตาลทุกถุง)
\concept{นี่คือการประมาณค่าเฉลี่ยของ \textbf{1 ประชากร} เพราะมีตัวแปรที่สนใจเพียงตัวเดียว (น้ำหนักน้ำตาล) และมีกลุ่มตัวอย่างเพียงกลุ่มเดียว}
\end{stepbox}

\begin{stepbox}{ขั้นที่ 2: กำหนดระดับความเชื่อมั่น}
โจทย์กำหนดความเชื่อมั่น 95\% $\Rightarrow 1-\alpha=0.95 \Rightarrow \alpha=0.05 \Rightarrow \alpha/2=0.025$
\end{stepbox}

\begin{stepbox}{ขั้นที่ 3: เลือกตัวสถิติที่ใช้ในการประมาณค่า}
เนื่องจาก\textbf{ทราบค่า $\sigma$ ของประชากรอยู่แล้ว} (ไม่ต้องสนใจว่า $n=25$ มากหรือน้อยกว่า 30) จึงตรงกับ\textbf{กรณีที่ 1} ใช้สูตร
$$\bar{X}\pm z_{\frac{\alpha}{2}}\dfrac{\sigma}{\sqrt{n}}$$
\concept{ตัวชี้วัดสำคัญที่สุดในการเลือกระหว่างสูตร $Z$ กับ $T$ คือ ``ทราบ $\sigma$ ของประชากรหรือไม่'' ถ้าทราบ ให้ใช้ $Z$ เสมอไม่ว่า $n$ จะเท่าใด โจทย์นี้ให้ $\sigma=1.2$ มาโดยตรง จึงไม่ต้องพิจารณาขนาดตัวอย่างเลย}
\end{stepbox}

\begin{stepbox}{ขั้นที่ 4: คำนวณ}
หาค่าวิกฤต $z_{\alpha/2}=z_{0.025}$ จากตาราง:
$$z_{0.025}=1.96$$
คำนวณค่าคลาดเคลื่อนมาตรฐาน (Standard Error):
$$\frac{\sigma}{\sqrt n}=\frac{1.2}{\sqrt{25}}=\frac{1.2}{5}=0.24$$
คำนวณค่าคลาดเคลื่อนสูงสุดที่ยอมรับได้ (Margin of Error):
$$z_{0.025}\times\frac{\sigma}{\sqrt n}=1.96\times0.24=0.4704$$
แทนค่าในสูตรช่วงความเชื่อมั่น:
$$19.8-0.4704 \;<\; \mu \;<\; 19.8+0.4704$$
$$19.3296 \;<\; \mu \;<\; 20.2704$$
ปัดเศษเป็น 2 ตำแหน่งทศนิยม:
$$19.33 \;<\; \mu \;<\; 20.27$$
\end{stepbox}

\begin{stepbox}{ขั้นที่ 5: สรุปคำตอบ}
ที่ระดับความเชื่อมั่น 95\% น้ำหนักสุทธิเฉลี่ยของน้ำตาลที่บรรจุถุงโดยเครื่องจักรเครื่องนี้ทั้งหมด มีค่าอยู่ระหว่าง \textbf{19.33 ถึง 20.27 ออนซ์}
\end{stepbox}

\begin{qbox}
ในการศึกษาปริมาณการใช้น้ำมันของรถบรรทุกรุ่นใหม่ยี่ห้อหนึ่ง รวม 10 คัน โดยบันทึกระยะทางที่รถแต่ละคันแล่นได้ (ไมล์/แกลลอน) ดังนี้: 18.6, 18.4, 19.2, 20.8, 19.4, 20.5, 18.7, 19.3, 19.6, 20.2 จงหาความเชื่อมั่น 99\% ของระยะทางโดยเฉลี่ยที่รถบรรทุกรุ่นใหม่ยี่ห้อนี้แล่นได้
\end{qbox}

\begin{stepbox}{ขั้นที่ 1: ทำความเข้าใจโจทย์}
โจทย์ให้ข้อมูลดิบจากตัวอย่างรถ $n=10$ คัน \textbf{ไม่ได้บอก $\sigma$ ของประชากรมาให้} ต้องคำนวณค่าประมาณจากตัวอย่างเอง สิ่งที่ต้องการคือช่วงความเชื่อมั่นของระยะทางเฉลี่ย $\mu$ ที่รถบรรทุกรุ่นนี้ทั้งหมด (ทุกคันที่ผลิต) แล่นได้
\concept{ยังคงเป็นการประมาณค่าเฉลี่ยของ 1 ประชากร แต่ต่างจากตัวอย่างก่อนหน้าตรงที่โจทย์นี้\textbf{ไม่ให้ $\sigma$ มา} จึงต้องประมาณ $\sigma$ ด้วย $S$ ที่คำนวณจากข้อมูลตัวอย่างเอง}
\end{stepbox}

\begin{stepbox}{ขั้นที่ 2: กำหนดระดับความเชื่อมั่น}
โจทย์กำหนดความเชื่อมั่น 99\% $\Rightarrow \alpha=0.01 \Rightarrow \alpha/2=0.005$
\end{stepbox}

\begin{stepbox}{ขั้นที่ 3: เลือกตัวสถิติที่ใช้ในการประมาณค่า}
ไม่ทราบ $\sigma$ ของประชากร และขนาดตัวอย่าง $n=10<30$ (ต้องสมมติว่าปริมาณน้ำมันที่ใช้แจกแจงแบบปกติ) จึงตรงกับ\textbf{กรณีที่ 2} ใช้สูตร
$$\bar{X}\pm t_{\frac{\alpha}{2}}\dfrac{S}{\sqrt{n}},\quad v=n-1$$
\concept{จุดสังเกต: ถ้า $n<30$ และไม่ทราบ $\sigma$ จะใช้ $T$-distribution เสมอ ไม่ใช่ $Z$ เพราะเมื่อขนาดตัวอย่างเล็ก การประมาณ $\sigma$ ด้วย $S$ มีความไม่แน่นอนสูง การแจกแจง $T$ จะมีหางที่หนากว่าปกติเพื่อชดเชยความไม่แน่นอนนี้}
\end{stepbox}

\begin{stepbox}{ขั้นที่ 4: คำนวณ}
คำนวณค่าเฉลี่ยตัวอย่าง:
$$\bar{X}=\frac{\sum X}{n}=\frac{18.6+18.4+\cdots+20.2}{10}=\frac{194.7}{10}=19.47 \text{ ไมล์/แกลลอน}$$
คำนวณผลรวมกำลังสอง $\sum X^2$ ของข้อมูลทั้ง 10 ค่า:
$$\sum X^2 = 18.6^2+18.4^2+\cdots+20.2^2 = 3{,}796.79$$
คำนวณความแปรปรวนตัวอย่าง ด้วยสูตรลัด:
$$S^2=\frac{\sum X^2-\dfrac{(\sum X)^2}{n}}{n-1}=\frac{3{,}796.79-\dfrac{(194.7)^2}{10}}{9}=\frac{3{,}796.79-3{,}790.809}{9}=\frac{5.981}{9}=0.66$$
$$S=\sqrt{0.66}=0.81 \text{ ไมล์/แกลลอน}$$
หาค่าวิกฤต $t_{\alpha/2}$ โดยองศาอิสระ $v=n-1=10-1=9$:
$$t_{0.005,\,9}=3.25$$
คำนวณค่าคลาดเคลื่อนสูงสุด:
$$t_{0.005}\times\frac{S}{\sqrt n}=3.25\times\frac{0.81}{\sqrt{10}}=3.25\times0.2562=0.8327$$
แทนค่าในสูตรช่วงความเชื่อมั่น:
$$19.47-0.8327 \;<\; \mu \;<\; 19.47+0.8327$$
$$18.64 \;<\; \mu \;<\; 20.30 \text{ ไมล์/แกลลอน}$$
\end{stepbox}

\begin{stepbox}{ขั้นที่ 5: สรุปคำตอบ}
ที่ระดับความเชื่อมั่น 99\% ระยะทางโดยเฉลี่ยที่รถบรรทุกรุ่นใหม่ยี่ห้อนี้แล่นได้ มีค่าอยู่ระหว่าง \textbf{18.64 ถึง 20.30 ไมล์/แกลลอน}
\end{stepbox}

\spssout{Excel (Data Analysis $\to$ Descriptive Statistics) จากไฟล์ต้นฉบับ}
คำสั่ง: \texttt{Data $\to$ Data Analysis $\to$ Descriptive Statistics} ระบุ Input Range เป็นช่วงข้อมูล 10 ค่า, ติ๊ก \texttt{Summary statistics} และ \texttt{Confidence Level for Mean} ใส่ 99\%
\begin{center}\small
\begin{tabular}{@{}lr@{}}
\toprule
\multicolumn{2}{c}{ระยะทาง (ไมล์/แกลลอน)}\\
\midrule
Mean & 19.47\\
Standard Error & 0.257789751\\
Median & 19.35\\
Mode & \#N/A\\
Standard Deviation & 0.81520277\\
Sample Variance & 0.664555556\\
Kurtosis & $-0.98854444$\\
Skewness & 0.377266757\\
Range & 2.4\\
Minimum & 18.4\\
Maximum & 20.8\\
Sum & 194.7\\
Count & 10\\
\textbf{Confidence Level(99.0\%)} & \textbf{0.837774294}\\
\bottomrule
\end{tabular}
\end{center}
\concept{Excel รายงาน Mean$=19.47$ ตรงกับค่าประมาณแบบจุดที่คำนวณด้วยมือ และรายงาน ``Confidence Level(99.0\%)$=0.8378$'' ซึ่งก็คือ\textbf{ค่าคลาดเคลื่อนสูงสุด} (margin of error, $t_{\alpha/2}\cdot S/\sqrt n$) ตัวเดียวกับที่คำนวณได้ข้างต้น ($0.8327\approx0.838$ ต่างกันเล็กน้อยจากการปัดเศษ) เมื่อนำไปบวกลบกับ Mean จะได้ช่วง $19.47\pm0.838 \Rightarrow 18.63<\mu<20.31$ ใกล้เคียงกับคำตอบที่คำนวณด้วยมือ}

\begin{qbox}
งานวิจัยเรื่อง ``ผลกระทบจากการแพร่ระบาดของโรคโควิด-19 ต่อความเครียดของนักศึกษาระดับปริญญาตรี'' สุ่มนักศึกษามา 150 คน ได้คะแนนความเครียดเฉลี่ย $\bar X=20.88$ และส่วนเบี่ยงเบนมาตรฐานตัวอย่าง $S=7.13$ จงประมาณค่าเฉลี่ยความเครียด ($\mu$) ของนักศึกษาทั้งหมดในประชากรที่ระดับความเชื่อมั่น 99\%
\end{qbox}

\begin{stepbox}{ขั้นที่ 1: ทำความเข้าใจโจทย์}
โจทย์ให้ $n=150$ (ตัวอย่างขนาดใหญ่), $\bar X=20.88$, $S=7.13$ และ\textbf{ไม่ทราบ $\sigma$ ของประชากร} สิ่งที่ต้องการคือช่วงความเชื่อมั่นของคะแนนความเครียดเฉลี่ย $\mu$ ของนักศึกษาทุกคนในประชากร (ไม่ใช่เฉพาะ 150 คนที่สุ่มมา)
\end{stepbox}

\begin{stepbox}{ขั้นที่ 2: กำหนดระดับความเชื่อมั่น}
โจทย์กำหนดความเชื่อมั่น 99\% $\Rightarrow \alpha=0.01 \Rightarrow \alpha/2=0.005$
\end{stepbox}

\begin{stepbox}{ขั้นที่ 3: เลือกตัวสถิติที่ใช้ในการประมาณค่า}
ไม่ทราบ $\sigma$ แต่ขนาดตัวอย่าง $n=150\ge30$ จึงตรงกับ\textbf{กรณีที่ 3} ใช้สูตร
$$\bar{X}\pm z_{\frac{\alpha}{2}}\dfrac{S}{\sqrt{n}}$$
\concept{แม้จะไม่ทราบ $\sigma$ เหมือนตัวอย่างรถบรรทุก แต่ครั้งนี้ $n=150\ge30$ (ตัวอย่างขนาดใหญ่) จึงใช้ $Z$ แทน $T$ ได้ โดยประมาณ $\sigma$ ด้วย $S$ อาศัยทฤษฎีบทลิมิตกลาง (Central Limit Theorem) ที่การแจกแจงของ $\bar X$ เข้าใกล้โค้งปกติเมื่อ $n$ มากพอ}
\end{stepbox}

\begin{stepbox}{ขั้นที่ 4: คำนวณ}
หาค่าวิกฤต:
$$z_{0.005}=2.576$$
คำนวณค่าคลาดเคลื่อนมาตรฐาน:
$$\frac{S}{\sqrt n}=\frac{7.13}{\sqrt{150}}=\frac{7.13}{12.247}=0.5822$$
คำนวณค่าคลาดเคลื่อนสูงสุด:
$$z_{0.005}\times\frac{S}{\sqrt n}=2.576\times0.5822=1.4997\approx1.50$$
แทนค่าในสูตรช่วงความเชื่อมั่น:
$$20.88-1.50 \;<\;\mu\;<\;20.88+1.50$$
$$19.38 \;<\; \mu \;<\; 22.38$$
\end{stepbox}

\begin{stepbox}{ขั้นที่ 5: สรุปคำตอบ}
ที่ระดับความเชื่อมั่น 99\% คะแนนความเครียดเฉลี่ยของนักศึกษาทั้งหมด (ประชากร) มีค่าอยู่ระหว่าง \textbf{19.38 ถึง 22.38 คะแนน}
\end{stepbox}

\spssout{SPSS (Analyze $\to$ Descriptive Statistics $\to$ Explore) จากไฟล์ต้นฉบับ}
คำสั่ง: \texttt{Analyze $\to$ Descriptive Statistics $\to$ Explore\dots} ส่งตัวแปรคะแนนความเครียดไปที่ \texttt{Dependent List} จากนั้นกดปุ่ม \texttt{Statistics} แล้วเปลี่ยน \texttt{Confidence Interval for Mean} เป็น 99
\begin{center}\small
\begin{tabular}{@{}llr@{}}
\toprule
\multicolumn{3}{l}{Descriptives (คะแนนความเครียด)}\\
\midrule
\multicolumn{2}{l}{Mean} & 20.88 (SE $=$ .582)\\
\multirow{2}{*}{99\% Confidence Interval for Mean} & Lower Bound & 19.36\\
 & Upper Bound & 22.40\\
\multicolumn{2}{l}{5\% Trimmed Mean} & 20.94\\
\multicolumn{2}{l}{Median} & 21.00\\
\multicolumn{2}{l}{Variance} & 50.777\\
\multicolumn{2}{l}{Std. Deviation} & 7.126\\
\multicolumn{2}{l}{Minimum / Maximum} & 1 / 39\\
\multicolumn{2}{l}{Range} & 38\\
\multicolumn{2}{l}{Interquartile Range} & 9\\
\multicolumn{2}{l}{Skewness / Kurtosis} & $-.116$ / .062\\
\bottomrule
\end{tabular}
\end{center}
\concept{การประมาณแบบจุด: $\bar X=20.88$ ตรงกับที่คำนวณด้วยมือ; การประมาณแบบช่วง 99\% CI: $19.36<\mu<22.40$ ใกล้เคียงกับคำตอบที่คำนวณด้วยมือ ($19.38<\mu<22.38$) ความต่างเล็กน้อยเกิดจาก SPSS ใช้ค่า $S=7.126$ เต็มจำนวนทศนิยมในการคำนวณ ในขณะที่การคำนวณด้วยมือปัดเศษ $S$ เป็น 7.13 ก่อน}

\section{ผลต่างค่าเฉลี่ย 2 ประชากร (ตัวอย่างอิสระกัน) $(\mu_1-\mu_2)$}

ใช้เมื่อต้องการเปรียบเทียบค่าเฉลี่ยของ\textbf{สองกลุ่มที่เป็นคนละกลุ่มกัน} (ไม่ใช่วัดจากคนกลุ่มเดียวกันซ้ำสองครั้ง) เช่น เปรียบเทียบคะแนนสอบระหว่างกลุ่มที่เรียนพิเศษกับกลุ่มที่ไม่เรียนพิเศษ

\dist{การประมาณแบบจุด}
\begin{formulabox}
$$\bar{X}_1-\bar{X}_2=\frac{\sum X_{1i}}{n_1}-\frac{\sum X_{2i}}{n_2}$$
\end{formulabox}

\dist{การประมาณแบบช่วง — 3 กรณี}
\begin{formulabox}
\case{กรณี 1 (ทราบ $\sigma_1^2,\sigma_2^2$ ของทั้งสองประชากร):}
$$(\bar{X}_1-\bar{X}_2)\pm z_{\frac{\alpha}{2}}\sqrt{\dfrac{\sigma_1^2}{n_1}+\dfrac{\sigma_2^2}{n_2}}$$
\case{กรณี 2 (ไม่ทราบ $\sigma_1^2,\sigma_2^2$ แต่ทราบว่า $\sigma_1^2=\sigma_2^2$, $n_1,n_2<30$):}
$$(\bar{X}_1-\bar{X}_2)\pm t_{\frac{\alpha}{2}}\sqrt{S_p^2\!\left(\dfrac1{n_1}+\dfrac1{n_2}\right)}$$
$$S_p^2=\dfrac{(n_1-1)S_1^2+(n_2-1)S_2^2}{n_1+n_2-2},\qquad v=n_1+n_2-2$$
\case{กรณี 3 (ไม่ทราบ $\sigma_1^2,\sigma_2^2$ และทราบว่า $\sigma_1^2\ne\sigma_2^2$, $n_1,n_2<30$):}
$$(\bar{X}_1-\bar{X}_2)\pm t_{\frac{\alpha}{2}}\sqrt{\dfrac{S_1^2}{n_1}+\dfrac{S_2^2}{n_2}}$$
$$v=\dfrac{\left(\dfrac{S_1^2}{n_1}+\dfrac{S_2^2}{n_2}\right)^2}{\dfrac{(S_1^2/n_1)^2}{n_1-1}+\dfrac{(S_2^2/n_2)^2}{n_2-1}} \quad\text{(สูตร Welch--Satterthwaite)}$$
\end{formulabox}

\textbf{ตัวแปรสำคัญ:} $S_p^2$ = ความแปรปรวนรวม (pooled variance) ใช้เมื่อสมมติว่าทั้งสองประชากรมีความแปรปรวนเท่ากัน (ผสานข้อมูลทั้งสองกลุ่มเข้าด้วยกันเพื่อประมาณ $\sigma^2$ ตัวเดียว); $v$ ในกรณี 3 มักได้ค่าไม่เป็นจำนวนเต็ม จึง\textbf{ปัดลง} เป็นจำนวนเต็มเสมอเพื่อความระมัดระวัง (ให้ช่วงกว้างกว่าเล็กน้อย)

\begin{qbox}
ในการศึกษาผลของการเรียนพิเศษหลังเลิกเรียน ต่อคะแนนสอบวิชาคณิตศาสตร์ของนักเรียนชั้นมัธยมศึกษาปีที่ 3 โรงเรียนแห่งหนึ่ง โดยสุ่มนักเรียนมา 2 กลุ่ม คือ กลุ่มที่เรียนพิเศษจำนวน 8 คน (148, 171, 198, 168, 218, 236, 178, 264) และกลุ่มที่ไม่ได้เรียนพิเศษจำนวน 24 คน (175, 172, 163, 181, 162, 152, 164, 180, 160, 174, 178, 148, 146, 176, 185, 158, 157, 164, 182, 170, 172, 178, 154, 148) เมื่อทราบว่า $\sigma_1^2\ne\sigma_2^2$ จงหาช่วงความเชื่อมั่น 90\% ของผลต่างของค่าเฉลี่ยคะแนนสอบระหว่างนักเรียนที่เรียนพิเศษและไม่ได้เรียนพิเศษ
\end{qbox}

\begin{stepbox}{ขั้นที่ 1: ทำความเข้าใจโจทย์}
มีตัวอย่าง\textbf{สองกลุ่มที่เป็นอิสระต่อกัน} (คนละกลุ่มคน ไม่ใช่วัดซ้ำ): กลุ่มเรียนพิเศษ $n_1=8$ และกลุ่มไม่เรียนพิเศษ $n_2=24$ โจทย์บอกมาแล้วว่า $\sigma_1^2\ne\sigma_2^2$ (ความแปรปรวนสองประชากรไม่เท่ากัน) สิ่งที่ต้องการคือช่วงความเชื่อมั่นของผลต่างค่าเฉลี่ย $\mu_1-\mu_2$
\concept{เป็นการประมาณผลต่างค่าเฉลี่ยของ \textbf{2 ประชากรอิสระกัน} เพราะนักเรียนสองกลุ่มเป็นคนละคนกัน ไม่ใช่คนกลุ่มเดียวกันที่ถูกวัดสองครั้ง}
\end{stepbox}

\begin{stepbox}{ขั้นที่ 2: กำหนดระดับความเชื่อมั่น}
โจทย์กำหนดความเชื่อมั่น 90\% $\Rightarrow \alpha=0.10 \Rightarrow \alpha/2=0.05$
\end{stepbox}

\begin{stepbox}{ขั้นที่ 3: เลือกตัวสถิติที่ใช้ในการประมาณค่า}
ไม่ทราบ $\sigma_1^2,\sigma_2^2$ และทราบว่า $\sigma_1^2\ne\sigma_2^2$ ทั้ง $n_1=8$ และ $n_2=24$ ต่างก็ $<30$ จึงตรงกับ\textbf{กรณีที่ 3} ใช้สูตร
$$(\bar{X}_1-\bar{X}_2)\pm t_{\frac{\alpha}{2}}\sqrt{\dfrac{S_1^2}{n_1}+\dfrac{S_2^2}{n_2}}, \qquad v=\dfrac{\left(\dfrac{S_1^2}{n_1}+\dfrac{S_2^2}{n_2}\right)^2}{\dfrac{(S_1^2/n_1)^2}{n_1-1}+\dfrac{(S_2^2/n_2)^2}{n_2-1}}$$
\concept{ข้อมูลที่โจทย์ ``บอกมาให้'' ว่า $\sigma_1^2\ne\sigma_2^2$ คือกุญแจสำคัญที่ทำให้เลือกกรณี 3 แทนกรณี 2 -- ถ้าโจทย์บอกว่าความแปรปรวนเท่ากันแทน จะต้องใช้ $S_p^2$ (pooled) ในกรณี 2 แทน}
\end{stepbox}

\begin{stepbox}{ขั้นที่ 4: คำนวณ}
คำนวณค่าเฉลี่ยแต่ละกลุ่ม:
$$\bar X_1=\frac{148+171+\cdots+264}{8}=\frac{1{,}581}{8}=197.625$$
$$\bar X_2=\frac{175+172+\cdots+148}{24}=\frac{4{,}035}{24}=168.125$$
ค่าประมาณแบบจุดของผลต่าง:
$$\bar X_1-\bar X_2=197.625-168.125=29.5$$
คำนวณความแปรปรวนตัวอย่างแต่ละกลุ่ม ด้วยสูตรลัด ($\sum X_1^2=323{,}173$ และ $\sum X_2^2=681{,}537$):
$$S_1^2=\frac{\sum X_1^2-\dfrac{(\sum X_1)^2}{n_1}}{n_1-1}=\frac{323{,}173-\dfrac{(1{,}581)^2}{8}}{7}=\frac{323{,}173-312{,}445.125}{7}=\frac{10{,}727.875}{7}=1{,}532.55$$
$$S_2^2=\frac{\sum X_2^2-\dfrac{(\sum X_2)^2}{n_2}}{n_2-1}=\frac{681{,}537-\dfrac{(4{,}035)^2}{24}}{23}=\frac{681{,}537-678{,}330.375}{23}=\frac{3{,}152.625}{23}=137.07$$
คำนวณองศาอิสระ $v$ ทีละส่วน เริ่มจากตัวตั้ง (ตัวเศษ):
$$\frac{S_1^2}{n_1}=\frac{1{,}532.55}{8}=191.57 \qquad \frac{S_2^2}{n_2}=\frac{137.07}{24}=5.71$$
$$\left(\frac{S_1^2}{n_1}+\frac{S_2^2}{n_2}\right)^2=(191.57+5.71)^2=(197.28)^2=38{,}919.4$$
คำนวณตัวหาร (ผลรวม 2 พจน์):
$$\frac{(S_1^2/n_1)^2}{n_1-1}=\frac{191.57^2}{7}=\frac{36{,}698.7}{7}=5{,}242.7 \qquad \frac{(S_2^2/n_2)^2}{n_2-1}=\frac{5.71^2}{23}=\frac{32.6}{23}=1.4$$
$$v=\frac{38{,}919.4}{5{,}242.7+1.4}=\frac{38{,}919.4}{5{,}244.1}=7.42 \approx 7 \text{ (ปัดลง)}$$
หาค่าวิกฤต $t_{0.05,\,7}=1.895$ แล้วคำนวณค่าคลาดเคลื่อนสูงสุด:
$$t_{0.05}\sqrt{\frac{S_1^2}{n_1}+\frac{S_2^2}{n_2}}=1.895\times\sqrt{191.57+5.71}=1.895\times\sqrt{197.28}=1.895\times14.045=26.62$$
แทนค่าในสูตรช่วงความเชื่อมั่น:
$$29.5-26.62 \;<\;\mu_1-\mu_2\;<\;29.5+26.62$$
$$2.88 \;<\; \mu_1-\mu_2 \;<\; 56.12$$
\end{stepbox}

\begin{stepbox}{ขั้นที่ 5: สรุปคำตอบ}
ที่ระดับความเชื่อมั่น 90\% ผลต่างของค่าเฉลี่ยคะแนนสอบระหว่างนักเรียนที่เรียนพิเศษและไม่ได้เรียนพิเศษ มีค่าอยู่ระหว่าง \textbf{2.88 ถึง 56.12 คะแนน} เนื่องจากช่วงนี้เป็นบวกทั้งช่วง (ไม่คร่อมศูนย์) จึงบ่งชี้ว่ากลุ่มที่เรียนพิเศษน่าจะมีคะแนนเฉลี่ยสูงกว่ากลุ่มที่ไม่ได้เรียนพิเศษ
\end{stepbox}

\begin{qbox}
งานวิจัยเรื่องผลกระทบจากโควิด-19 ต่อความเครียดของนักศึกษา ($n=150$) อยากทราบว่านักศึกษาชายและนักศึกษาหญิงมีคะแนนความเครียดเฉลี่ยต่างกันเพียงใด จงประมาณค่าช่วงความเชื่อมั่น 95\% ของ $\mu_{\text{ชาย}}-\mu_{\text{หญิง}}$
\end{qbox}

\begin{stepbox}{ขั้นที่ 1: ทำความเข้าใจโจทย์}
เปรียบเทียบคะแนนความเครียดเฉลี่ยระหว่าง\textbf{นักศึกษาชาย}กับ\textbf{นักศึกษาหญิง} ซึ่งเป็นคนละกลุ่มกัน (ตัวอย่างอิสระ) ต้องการช่วงความเชื่อมั่นของ $\mu_1-\mu_2$ (ชาย$-$หญิง)
\end{stepbox}

\begin{stepbox}{ขั้นที่ 2: กำหนดระดับความเชื่อมั่น}
กำหนดความเชื่อมั่น 95\% $\Rightarrow \alpha=0.05$
\end{stepbox}

\begin{stepbox}{ขั้นที่ 3: เลือกตัวสถิติที่ใช้ในการประมาณค่า}
ไม่ทราบ $\sigma_1^2,\sigma_2^2$ ของประชากร จึงต้องตรวจสอบก่อนว่าความแปรปรวนสองกลุ่มเท่ากันหรือไม่ ด้วย\textbf{Levene's Test}:
$$H_0:\sigma_1^2=\sigma_2^2 \qquad H_1:\sigma_1^2\ne\sigma_2^2$$
\concept{ถ้า Sig. ของ Levene's Test $>\alpha$ แปลว่ายอมรับว่าความแปรปรวนเท่ากัน ให้อ่านผลจากแถว ``Equal variances assumed'' (สูตรกรณี 2, ใช้ $S_p^2$); ถ้า Sig. $\le\alpha$ ให้อ่านแถว ``Equal variances not assumed'' (สูตรกรณี 3, Welch) ขั้นตอนนี้ SPSS คำนวณให้อัตโนมัติในตารางเดียวกับผลลัพธ์หลัก}
\end{stepbox}

\begin{stepbox}{ขั้นที่ 4: คำนวณ}
คำสั่งใน SPSS: \texttt{Analyze $\to$ Compare Means $\to$ Independent-Samples T Test\dots} ส่งตัวแปรคะแนนความเครียดไปที่ \texttt{Test Variable(s)} และตัวแปรเพศไปที่ \texttt{Grouping Variable}

\spssout{SPSS จากไฟล์ต้นฉบับ}
\begin{center}\small
\begin{tabular}{@{}llcccc@{}}
\toprule
& เพศ & N & Mean & Std. Deviation & Std. Error Mean\\
\midrule
\multirow{2}{*}{คะแนนความเครียด} & ชาย & 106 & 21.87 & 6.991 & .679\\
& หญิง & 44 & 18.50 & 6.957 & 1.049\\
\bottomrule
\end{tabular}
\end{center}
\begin{center}\tiny
\begin{tabular}{@{}l cc ccc cc cc@{}}
\toprule
& \multicolumn{2}{c}{Levene's Test} & \multicolumn{7}{c}{t-test for Equality of Means}\\
\cmidrule(lr){2-3}\cmidrule(lr){4-10}
& F & Sig. & $t$ & $df$ & Sig.(2-t.) & Mean Diff. & SE Diff. & \multicolumn{2}{c}{95\% CI}\\
& & & & & & & & Lower & Upper\\
\midrule
Equal variances assumed & .033 & .856 & 2.690 & 148 & .008 & 3.368 & 1.252 & .894 & 5.842\\
Equal variances not assumed & & & 2.696 & 80.796 & .009 & 3.368 & 1.249 & .882 & 5.854\\
\bottomrule
\end{tabular}
\end{center}
ค่าประมาณแบบจุด: $\bar X_1-\bar X_2 = 21.87-18.50 = 3.368$

อ่านผล Levene's Test ก่อน: Sig.$=.856>0.05 \Rightarrow$ ยอมรับว่าความแปรปรวนสองกลุ่มเท่ากัน $\Rightarrow$ ใช้แถว \textbf{``Equal variances assumed''}:
$$0.894 \;<\; \mu_1-\mu_2 \;<\; 5.842$$
(หากใช้แถว ``Equal variances not assumed'' แทน จะได้ $0.882<\mu_1-\mu_2<5.854$ ซึ่งใกล้เคียงกันมากในกรณีนี้)
\end{stepbox}

\begin{stepbox}{ขั้นที่ 5: สรุปคำตอบ}
ที่ระดับความเชื่อมั่น 95\% ผลต่างของคะแนนความเครียดเฉลี่ยระหว่างนักศึกษาชายและหญิง (ชาย$-$หญิง) มีค่าอยู่ระหว่าง \textbf{0.894 ถึง 5.842 คะแนน} เนื่องจากช่วงนี้เป็นบวกทั้งช่วง (ไม่คร่อมศูนย์) จึงบ่งชี้ว่านักศึกษาชายมีแนวโน้มคะแนนความเครียดเฉลี่ยสูงกว่านักศึกษาหญิง
\end{stepbox}

\section{ผลต่างค่าเฉลี่ย 2 ประชากร (ตัวอย่างมีความสัมพันธ์กัน) $(\mu_D)$}

ใช้เมื่อข้อมูล 2 ชุดมาจาก\textbf{กลุ่มตัวอย่างเดียวกัน} ที่ถูกวัดสองครั้ง (เช่น วัดก่อน-หลังการทดลอง) ต่างจากหัวข้อที่แล้วซึ่งเป็นคนละกลุ่มกัน ให้ $D_i=X_{i1}-X_{i2}$, $i=1,2,\dots,n$ เป็นผลต่างของตัวอย่างสุ่มจากประชากรที่มีการแจกแจงแบบปกติ มีค่าเฉลี่ย $\mu_D$ และความแปรปรวน $\sigma_D^2$ โดยที่ $\mu_D=\mu_1-\mu_2$

ให้ $D_0$ เป็นค่าผลต่างของค่าเฉลี่ยที่สนใจทดสอบ: ถ้า $D_0=0$ หมายความว่าก่อนและหลังไม่ต่างกัน, ถ้า $D_0>0$ หมายความว่าก่อนทดลองมีค่าเฉลี่ยสูงกว่าหลังทดลอง, ถ้า $D_0<0$ หมายความว่าก่อนทดลองมีค่าเฉลี่ยต่ำกว่าหลังทดลอง

\begin{formulabox}
$$\bar D=\frac{\sum D_i}{n},\qquad S_D^2=\frac{\sum D_i^2-\dfrac{(\sum D_i)^2}{n}}{n-1}$$
\case{กรณี 1 ($n<30$):} $\displaystyle \bar D\pm t_{\frac{\alpha}{2}}\dfrac{S_D}{\sqrt n},\ v=n-1$\\[4pt]
\case{กรณี 2 ($n\ge30$):} $\displaystyle \bar D\pm z_{\frac{\alpha}{2}}\dfrac{S_D}{\sqrt n}$
\end{formulabox}

\begin{qbox}
ผู้วิจัยจัดโปรแกรมบำบัดความเครียดด้วย ``ดนตรีบำบัด'' ให้นักศึกษากลุ่มตัวอย่าง แล้ววัดคะแนนความเครียดก่อนและหลังเข้าร่วมโปรแกรม จากนักศึกษา 9 คน ได้ผลดังนี้

\begin{center}
\small
\begin{tabular}{c|ccccccccc}
คนที่ & 1&2&3&4&5&6&7&8&9\\\hline
ก่อน & 22&21&25&24&24&19&22&20&30\\
หลัง & 23&15&20&20&23&16&18&16&32\\
$D=$ก่อน$-$หลัง & $-1$&6&5&4&1&3&4&4&$-2$\\
\end{tabular}
\end{center}
จงหาช่วงความเชื่อมั่น 95\% ของผลต่างค่าเฉลี่ยความเครียดก่อนและหลังเข้าร่วมโปรแกรม ($\mu_D$)
\end{qbox}

\begin{stepbox}{ขั้นที่ 1: ทำความเข้าใจโจทย์}
ข้อมูลชุดนี้เป็นคะแนนความเครียด ``ก่อน-หลัง'' ที่วัดจาก\textbf{นักศึกษากลุ่มเดียวกัน} 9 คน (วัดซ้ำ) จึงเป็นกรณี\textbf{ตัวอย่างมีความสัมพันธ์กัน} ไม่ใช่กรณีอิสระ ต้องการช่วงความเชื่อมั่นของ $\mu_D=\mu_{\text{ก่อน}}-\mu_{\text{หลัง}}$
\concept{ข้อสังเกต: จำนวนข้อมูลคือ $n=9$ คู่ (ไม่ใช่ $9\times2=18$) เพราะแต่ละคู่ ``ก่อน-หลัง'' ของคนคนเดียวกันถูกหลอมรวมเป็นค่า $D_i$ เพียงค่าเดียวต่อคนก่อนนำไปคำนวณ}
\end{stepbox}

\begin{stepbox}{ขั้นที่ 2: กำหนดระดับความเชื่อมั่น}
กำหนดความเชื่อมั่น 95\% $\Rightarrow \alpha=0.05 \Rightarrow \alpha/2=0.025$
\end{stepbox}

\begin{stepbox}{ขั้นที่ 3: เลือกตัวสถิติที่ใช้ในการประมาณค่า}
ขนาดตัวอย่าง $n=9<30$ (สมมติว่าผลต่าง $D_i$ แจกแจงแบบปกติ) จึงตรงกับ\textbf{กรณีที่ 1} ใช้สูตร
$$\bar D\pm t_{\frac{\alpha}{2}}\dfrac{S_D}{\sqrt n},\qquad v=n-1$$
\concept{หลักการเลือกสูตรเหมือนกรณีค่าเฉลี่ย 1 ประชากรทุกประการ เพียงแต่นำ ``ผลต่าง'' $D_i$ มาแทนตัวแปรเดี่ยว $X_i$: ถ้า $n<30$ ใช้ $T$ ($v=n-1$), ถ้า $n\ge30$ ใช้ $Z$}
\end{stepbox}

\begin{stepbox}{ขั้นที่ 4: คำนวณ}
คำนวณผลรวมและผลรวมกำลังสองของ $D_i$:
$$\sum D_i = (-1)+6+5+4+1+3+4+4+(-2)=24$$
$$\sum D_i^2 = 1+36+25+16+1+9+16+16+4=124$$
คำนวณค่าเฉลี่ยของผลต่าง:
$$\bar D=\frac{\sum D_i}{n}=\frac{24}{9}=2.667$$
คำนวณความแปรปรวนของผลต่าง ด้วยสูตรลัด:
$$S_D^2=\frac{\sum D_i^2-\dfrac{(\sum D_i)^2}{n}}{n-1}=\frac{124-\dfrac{24^2}{9}}{8}=\frac{124-64}{8}=\frac{60}{8}=7.5$$
$$S_D=\sqrt{7.5}=2.739$$
หาค่าวิกฤต โดยองศาอิสระ $v=n-1=8$:
$$t_{0.025,\,8}=2.306$$
คำนวณค่าคลาดเคลื่อนสูงสุด:
$$t_{0.025}\times\frac{S_D}{\sqrt n}=2.306\times\frac{2.739}{\sqrt 9}=2.306\times\frac{2.739}{3}=2.306\times0.913=2.105$$
แทนค่าในสูตรช่วงความเชื่อมั่น:
$$2.667-2.105 \;<\;\mu_D\;<\;2.667+2.105$$
$$0.56 \;<\; \mu_D \;<\; 4.77$$
\end{stepbox}

\begin{stepbox}{ขั้นที่ 5: สรุปคำตอบ}
ที่ระดับความเชื่อมั่น 95\% ผลต่างของคะแนนความเครียดเฉลี่ยก่อนและหลังเข้าร่วมโปรแกรมดนตรีบำบัด มีค่าอยู่ระหว่าง \textbf{0.56 ถึง 4.77 คะแนน} เนื่องจากช่วงทั้งหมดเป็นบวก (ไม่คร่อมศูนย์ และ $D_0=0$ ไม่อยู่ในช่วง) จึงบ่งชี้ว่าคะแนนความเครียด\textbf{ก่อน}เข้าร่วมโปรแกรมสูงกว่า\textbf{หลัง}เข้าร่วมโปรแกรมจริง กล่าวคือ โปรแกรมดนตรีบำบัดมีแนวโน้มช่วยลดความเครียดได้
\end{stepbox}

\begin{outbox}
\textbf{หมายเหตุเกี่ยวกับข้อมูลชุดนี้:} ไฟล์ประกอบการสอน CH4-การประมาณค่า.pdf แสดงเฉพาะข้อมูลดิบของนักศึกษา 9 คนข้างต้น (ไม่มีภาพหน้าจอ SPSS ประกอบตัวอย่างนี้โดยตรง) ผลการคำนวณด้วยมือข้างต้นจึงเป็นคำตอบที่ตรงกับข้อมูลชุดนี้ทั้งหมด

ส่วนงานวิจัยเรื่องเดียวกัน (``ดนตรีบำบัดกับความเครียด'') เมื่อวิเคราะห์จากข้อมูลกลุ่มตัวอย่างเต็มจำนวน $n=150$ คน (ปรากฏใน SPSS Paired-Samples T Test ของบทที่ 5) ได้ผลดังนี้:
\begin{center}\small
\begin{tabular}{@{}lccccc@{}}
\toprule
Pair & Mean (ก่อน$-$หลัง) & SD & SE Mean & \multicolumn{2}{c}{95\% CI ของ $\mu_D$}\\
& & & & Lower & Upper\\
\midrule
ก่อน $-$ หลัง & 2.647 & 6.779 & .554 & 1.553 & 3.740\\
\bottomrule
\end{tabular}
\end{center}
ทิศทางของผลลัพธ์ (ก่อน$>$หลัง, ช่วงเป็นบวกทั้งช่วง) สอดคล้องกับข้อสรุปจากตัวอย่าง 9 คนข้างต้น เพียงแต่ช่วงจากข้อมูลเต็ม $n=150$ แคบกว่ามาก เพราะขนาดตัวอย่างใหญ่กว่าทำให้ค่าคลาดเคลื่อนมาตรฐานเล็กลง
\end{outbox}

\section{ประมาณค่าสัดส่วนของประชากร ($p$)}

ใช้เมื่อตัวแปรที่สนใจเป็นข้อมูลเชิงกลุ่ม (มี/ไม่มี, ใช่/ไม่ใช่) ไม่ใช่ข้อมูลต่อเนื่อง

\begin{center}
\begin{tikzpicture}
\node[font=\scriptsize] at (0,0) {ประชากร $N$: กลุ่มที่สนใจมีจำนวน $X$, $p=X/N$};
\node[font=\scriptsize] at (0,-0.5) {$\downarrow$ สุ่มตัวอย่าง};
\node[font=\scriptsize] at (0,-1) {ตัวอย่าง $n$: กลุ่มที่สนใจมีจำนวน $x$, $\hat p=x/n$};
\end{tikzpicture}
\end{center}

\begin{formulabox}
\textbf{แบบจุด:}\quad $\hat{p}=\dfrac{x}{n}$\quad ($x$ = จำนวนหน่วยตัวอย่างที่มีลักษณะที่สนใจ, $n$ = ขนาดตัวอย่างทั้งหมด)\\[3pt]
\textbf{แบบช่วง:}\quad
$$\hat{p}\pm z_{\frac{\alpha}{2}}\sqrt{\frac{\hat{p}\hat{q}}{n}}, \qquad \hat q=1-\hat p$$
\end{formulabox}

\begin{qbox}
สอบถามแม่บ้านในจังหวัดขอนแก่นจำนวน 400 คน ว่าเป็นโรคความดันโลหิตสูงหรือไม่ พบว่ามีแม่บ้าน 224 คนเป็นโรคความดันโลหิตสูง จงหาความเชื่อมั่น 90\% ของสัดส่วนแม่บ้านทั้งหมดในจังหวัดขอนแก่นที่เป็นโรคนี้
\end{qbox}

\begin{stepbox}{ขั้นที่ 1: ทำความเข้าใจโจทย์}
ตัวแปรที่สนใจคือ ``เป็นโรคความดันโลหิตสูงหรือไม่'' ซึ่งเป็นข้อมูลเชิงกลุ่ม (สัดส่วน) ไม่ใช่ค่าเฉลี่ย จากตัวอย่าง $n=400$ คน พบผู้เป็นโรค $x=224$ คน ต้องการช่วงความเชื่อมั่นของสัดส่วน $p$ ในประชากร (แม่บ้านทั้งหมดในจังหวัดขอนแก่น)
\end{stepbox}

\begin{stepbox}{ขั้นที่ 2: กำหนดระดับความเชื่อมั่น}
โจทย์กำหนดความเชื่อมั่น 90\% $\Rightarrow \alpha=0.10 \Rightarrow \alpha/2=0.05$
\end{stepbox}

\begin{stepbox}{ขั้นที่ 3: เลือกตัวสถิติที่ใช้ในการประมาณค่า}
เมื่อประมาณค่าสัดส่วนของ 1 ประชากร มีสูตรเดียว (อาศัยการประมาณด้วยโค้งปกติ เมื่อ $n$ มากพอ):
$$\hat{p}\pm z_{\frac{\alpha}{2}}\sqrt{\frac{\hat{p}\hat{q}}{n}}$$
\concept{ต่างจากกรณีค่าเฉลี่ยที่มี 3 กรณีให้เลือก (ทราบ/ไม่ทราบ $\sigma$, $n$ มาก/น้อย) การประมาณค่าสัดส่วนใช้สูตร $Z$ เพียงสูตรเดียวเสมอ เพราะ $\hat p\hat q/n$ เป็นทั้งค่าประมาณของความแปรปรวนและคำนวณได้จากข้อมูลตัวอย่างโดยตรงอยู่แล้ว}
\end{stepbox}

\begin{stepbox}{ขั้นที่ 4: คำนวณ}
คำนวณค่าประมาณแบบจุด:
$$\hat p=\frac{x}{n}=\frac{224}{400}=0.56 \qquad \hat q=1-\hat p=1-0.56=0.44$$
หาค่าวิกฤต:
$$z_{0.05}=1.645$$
คำนวณค่าคลาดเคลื่อนมาตรฐาน:
$$\sqrt{\frac{\hat p\hat q}{n}}=\sqrt{\frac{(0.56)(0.44)}{400}}=\sqrt{\frac{0.2464}{400}}=\sqrt{0.000616}=0.0248$$
คำนวณค่าคลาดเคลื่อนสูงสุด:
$$z_{0.05}\times\sqrt{\frac{\hat p\hat q}{n}}=1.645\times0.0248=0.0408$$
แทนค่าในสูตรช่วงความเชื่อมั่น:
$$0.56-0.0408 \;<\;p\;<\;0.56+0.0408$$
$$0.52 \;<\; p \;<\; 0.60$$
\end{stepbox}

\begin{stepbox}{ขั้นที่ 5: สรุปคำตอบ}
ที่ระดับความเชื่อมั่น 90\% สัดส่วนของแม่บ้านทั้งหมดในจังหวัดขอนแก่นที่เป็นโรคความดันโลหิตสูง มีค่าอยู่ระหว่าง \textbf{0.52 ถึง 0.60} (หรือ 52\% ถึง 60\%)
\end{stepbox}

\begin{qbox}
เกณฑ์ประเมินความเครียดกำหนดว่า นักศึกษาจะถูกประเมินว่ามีความเครียดเมื่อคะแนนความเครียด $\ge20$ คะแนน จากกลุ่มตัวอย่าง 150 คน พบว่ามีนักศึกษาที่ถูกประเมินว่ามีความเครียด 101 คน จงหาช่วงความเชื่อมั่น 95\% ของสัดส่วนนักศึกษาที่มีความเครียดในประชากร
\end{qbox}

\begin{stepbox}{ขั้นที่ 1: ทำความเข้าใจโจทย์}
ตัวแปร ``มีความเครียดหรือไม่'' (คะแนน $\ge20$ หรือไม่) เป็นข้อมูลเชิงกลุ่ม จากตัวอย่าง $n=150$ คน พบผู้มีความเครียด $x=101$ คน ต้องการช่วงความเชื่อมั่นของสัดส่วน $p$ ของนักศึกษาที่มีความเครียดในประชากรทั้งหมด
\end{stepbox}

\begin{stepbox}{ขั้นที่ 2: กำหนดระดับความเชื่อมั่น}
กำหนดความเชื่อมั่น 95\% $\Rightarrow \alpha=0.05 \Rightarrow \alpha/2=0.025$
\end{stepbox}

\begin{stepbox}{ขั้นที่ 3: เลือกตัวสถิติที่ใช้ในการประมาณค่า}
ใช้สูตรประมาณค่าสัดส่วน 1 ประชากร:
$$\hat{p}\pm z_{\frac{\alpha}{2}}\sqrt{\frac{\hat{p}\hat{q}}{n}}$$
\end{stepbox}

\begin{stepbox}{ขั้นที่ 4: คำนวณ}
คำสั่งใน SPSS: \texttt{Analyze $\to$ Compare Means and Proportions $\to$ One-Sample Proportions\dots} ส่งตัวแปรความเครียดไปที่ \texttt{Test Variable(s)} และกำหนด \texttt{Define Success} เป็นค่า 1 (มีความเครียด)

\spssout{SPSS จากไฟล์ต้นฉบับ}
\begin{center}\small
\begin{tabular}{@{}lccccc@{}}
\toprule
& Successes & Trials & Proportion & Asym. SE & 95\% CI\\
\midrule
มีความเครียด (Wald) & 101 & 150 & .673 & .038 & (.598, .748)\\
\bottomrule
\end{tabular}
\end{center}
ตรวจสอบด้วยการคำนวณมือ: $\hat p=101/150=0.673$, $\hat q=0.327$, $z_{0.025}=1.96$
$$\sqrt{\frac{(0.673)(0.327)}{150}}=\sqrt{0.001467}=0.0383 \qquad 1.96\times0.0383=0.0750$$
$$0.673-0.0750 \;<\;p\;<\;0.673+0.0750 \;\Rightarrow\; 0.598 \;<\;p\;<\;0.748$$
ตรงกับผลลัพธ์จาก SPSS พอดี
\end{stepbox}

\begin{stepbox}{ขั้นที่ 5: สรุปคำตอบ}
ที่ระดับความเชื่อมั่น 95\% สัดส่วนของนักศึกษาที่มีความเครียดในประชากร มีค่าอยู่ระหว่าง \textbf{0.598 ถึง 0.748} (หรือประมาณ 60\% ถึง 75\%)
\end{stepbox}

\section{ผลต่างสัดส่วนของ 2 ประชากร $(p_1-p_2)$}

\begin{formulabox}
\textbf{แบบจุด:}\quad $\hat p_1-\hat p_2=\dfrac{x_1}{n_1}-\dfrac{x_2}{n_2}$\\[3pt]
\textbf{แบบช่วง:}
$$(\hat p_1-\hat p_2)\pm z_{\frac{\alpha}{2}}\sqrt{\frac{\hat p_1\hat q_1}{n_1}+\frac{\hat p_2\hat q_2}{n_2}}$$
\end{formulabox}

\concept{ข้อสังเกตสำคัญ: ในการ\textbf{ประมาณค่า}ผลต่างสัดส่วน (บทนี้) ให้ใช้ $\hat p_1,\hat p_2$ แยกกันคำนวณค่าคลาดเคลื่อนมาตรฐานเสมอ \textbf{ไม่ต้อง pool} รวมกันเป็น $\hat p$ ตัวเดียวแบบที่ทำในการทดสอบสมมติฐาน (บทที่ 5) เพราะการ pool มีเหตุผลมาจากการตั้งสมมติฐานว่า $p_1=p_2$ ซึ่งเป็นแนวคิดของการทดสอบสมมติฐาน ไม่ใช่การประมาณค่า}

\begin{qbox}
สุ่มวัยรุ่นหญิง 600 คน และวัยรุ่นชาย 400 คน ในจังหวัดหนึ่ง แล้วสอบถามว่าชมข่าวโทรทัศน์ภาคค่ำเป็นประจำหรือไม่ ปรากฏว่ามีวัยรุ่นหญิง 300 คน และวัยรุ่นชาย 100 คน ชมรายการดังกล่าวเป็นประจำ จงหาช่วงความเชื่อมั่น 99\% ของผลต่างสัดส่วนวัยรุ่นเพศหญิงและเพศชายทั้งหมดในจังหวัดนี้ที่ชมข่าวโทรทัศน์ภาคค่ำเป็นประจำ ($p_1-p_2$ โดย $p_1$ แทนหญิง, $p_2$ แทนชาย)
\end{qbox}

\begin{stepbox}{ขั้นที่ 1: ทำความเข้าใจโจทย์}
เปรียบเทียบ\textbf{สัดส่วน}ของสองกลุ่มอิสระ (หญิง $n_1=600$ พบ $x_1=300$ คน; ชาย $n_2=400$ พบ $x_2=100$ คน) ต้องการช่วงความเชื่อมั่นของผลต่างสัดส่วน $p_1-p_2$
\end{stepbox}

\begin{stepbox}{ขั้นที่ 2: กำหนดระดับความเชื่อมั่น}
โจทย์กำหนดความเชื่อมั่น 99\% $\Rightarrow \alpha=0.01 \Rightarrow \alpha/2=0.005$
\end{stepbox}

\begin{stepbox}{ขั้นที่ 3: เลือกตัวสถิติที่ใช้ในการประมาณค่า}
ใช้สูตรประมาณผลต่างสัดส่วน 2 ประชากร (ไม่ pool):
$$(\hat p_1-\hat p_2)\pm z_{\frac{\alpha}{2}}\sqrt{\frac{\hat p_1\hat q_1}{n_1}+\frac{\hat p_2\hat q_2}{n_2}}$$
\end{stepbox}

\begin{stepbox}{ขั้นที่ 4: คำนวณ}
คำนวณสัดส่วนตัวอย่างแต่ละกลุ่ม:
$$\hat p_1=\frac{300}{600}=0.5,\quad \hat q_1=0.5 \qquad\qquad \hat p_2=\frac{100}{400}=0.25,\quad \hat q_2=0.75$$
ค่าประมาณแบบจุดของผลต่าง:
$$\hat p_1-\hat p_2=0.5-0.25=0.25$$
หาค่าวิกฤต:
$$z_{0.005}=2.576$$
คำนวณค่าคลาดเคลื่อนมาตรฐาน ทีละพจน์:
$$\frac{\hat p_1\hat q_1}{n_1}=\frac{(0.5)(0.5)}{600}=0.000417 \qquad \frac{\hat p_2\hat q_2}{n_2}=\frac{(0.25)(0.75)}{400}=0.000469$$
$$\sqrt{\frac{\hat p_1\hat q_1}{n_1}+\frac{\hat p_2\hat q_2}{n_2}}=\sqrt{0.000417+0.000469}=\sqrt{0.000885}=0.02976$$
คำนวณค่าคลาดเคลื่อนสูงสุด:
$$z_{0.005}\times0.02976=2.576\times0.02976=0.0767$$
แทนค่าในสูตรช่วงความเชื่อมั่น:
$$0.25-0.0767\;<\;p_1-p_2\;<\;0.25+0.0767$$
$$0.17 \;<\; p_1-p_2 \;<\; 0.33$$
\end{stepbox}

\begin{stepbox}{ขั้นที่ 5: สรุปคำตอบ}
ที่ระดับความเชื่อมั่น 99\% ผลต่างของสัดส่วนวัยรุ่นเพศหญิงและเพศชายที่ชมข่าวโทรทัศน์ภาคค่ำเป็นประจำ (หญิง$-$ชาย) มีค่าอยู่ระหว่าง \textbf{0.17 ถึง 0.33} เนื่องจากช่วงเป็นบวกทั้งช่วง (ไม่คร่อมศูนย์) จึงบ่งชี้ว่าวัยรุ่นหญิงมีสัดส่วนการชมข่าวภาคค่ำเป็นประจำสูงกว่าวัยรุ่นชาย
\end{stepbox}

\begin{qbox}
จากงานวิจัยผลกระทบโควิด-19 ต่อความเครียดนักศึกษา อยากทราบว่านักศึกษาชายกับนักศึกษาหญิงมีสัดส่วนความเครียดต่างกันเพียงใด ข้อมูล: ชาย $x_1=74$ คนจาก $n_1=106$ คน; หญิง $x_2=27$ คนจาก $n_2=44$ คน จงหาช่วงความเชื่อมั่น 95\% ของ $p_1-p_2$ (ชาย$-$หญิง)
\end{qbox}

\begin{stepbox}{ขั้นที่ 1: ทำความเข้าใจโจทย์}
เปรียบเทียบสัดส่วนความเครียดระหว่างนักศึกษาชายและหญิง (สองกลุ่มอิสระ) ต้องการช่วงความเชื่อมั่นของ $p_1-p_2$
\end{stepbox}

\begin{stepbox}{ขั้นที่ 2: กำหนดระดับความเชื่อมั่น}
กำหนดความเชื่อมั่น 95\% $\Rightarrow \alpha=0.05 \Rightarrow \alpha/2=0.025$
\end{stepbox}

\begin{stepbox}{ขั้นที่ 3: เลือกตัวสถิติที่ใช้ในการประมาณค่า}
ใช้สูตรประมาณผลต่างสัดส่วน 2 ประชากร (ไม่ pool):
$$(\hat p_1-\hat p_2)\pm z_{\frac{\alpha}{2}}\sqrt{\frac{\hat p_1\hat q_1}{n_1}+\frac{\hat p_2\hat q_2}{n_2}}$$
\end{stepbox}

\begin{stepbox}{ขั้นที่ 4: คำนวณ}
คำสั่งใน SPSS: \texttt{Analyze $\to$ Compare Means and Proportions $\to$ Independent-Sample Proportions\dots} ส่งตัวแปรความเครียดไปที่ \texttt{Test Variable(s)} และตัวแปรเพศไปที่ \texttt{Grouping Variable}

\spssout{SPSS จากไฟล์ต้นฉบับ}
\begin{center}\small
\begin{tabular}{@{}lcccc@{}}
\toprule
เพศ & Successes & Trials & Proportion & Asym. SE\\
\midrule
ชาย & 74 & 106 & .698 & .045\\
หญิง & 27 & 44 & .614 & .073\\
\bottomrule
\end{tabular}

\begin{tabular}{@{}lcccc@{}}
\toprule
& Diff. in Proportions & Asym. SE & \multicolumn{2}{c}{95\% CI}\\
& & & Lower & Upper\\
\midrule
มีความเครียด (Wald) & .084 & .086 & $-.084$ & .253\\
\bottomrule
\end{tabular}
\end{center}
ค่าประมาณแบบจุด: $\hat p_1-\hat p_2=0.698-0.614=0.084$

ช่วงความเชื่อมั่น 95\%: $-0.084 \;<\; p_1-p_2 \;<\; 0.253$
\end{stepbox}

\begin{stepbox}{ขั้นที่ 5: สรุปคำตอบ}
ที่ระดับความเชื่อมั่น 95\% ผลต่างของสัดส่วนความเครียดระหว่างนักศึกษาชายและหญิง มีค่าอยู่ระหว่าง \textbf{$-0.084$ ถึง 0.253} เนื่องจากช่วงนี้\textbf{คร่อมค่า 0} (มีทั้งค่าบวกและค่าลบ) จึงยังสรุปทิศทางที่ชัดเจนไม่ได้ว่าเพศใดมีสัดส่วนความเครียดสูงกว่ากัน -- เป็นไปได้ว่าสัดส่วนความเครียดของชายและหญิงในประชากรอาจไม่แตกต่างกันเลยก็ได้
\end{stepbox}

\section{สรุปสูตรตัวประมาณแบบจุด}
\begin{formulabox}
\begin{tabular}{@{}ll@{}}
\textbf{พารามิเตอร์} & \textbf{ตัวประมาณแบบจุด}\\\hline
$\mu$ & $\bar X$\\
$\mu_1-\mu_2$ (อิสระ) & $\bar X_1-\bar X_2$\\
$\mu_D$ (สัมพันธ์กัน) & $\bar D$\\
$p$ & $\hat p=x/n$\\
$p_1-p_2$ & $\hat p_1-\hat p_2$\\
\end{tabular}
\end{formulabox}

\section{หลักการเลือกสูตรช่วงความเชื่อมั่น (สรุปรวม)}
\begin{center}\small
\begin{tabular}{@{}p{6.7cm}p{6.3cm}@{}}
\toprule
\textbf{สถานการณ์} & \textbf{สูตรที่ใช้}\\\midrule
$\mu$ (1 ประชากร), ทราบ $\sigma^2$ & $Z$ ($n$ เท่าใดก็ได้)\\
$\mu$ (1 ประชากร), ไม่ทราบ $\sigma^2$, $n<30$ & $T,\ v=n-1$\\
$\mu$ (1 ประชากร), ไม่ทราบ $\sigma^2$, $n\ge30$ & $Z$ (ประมาณด้วย $S$)\\
$\mu_1-\mu_2$ อิสระกัน, ทราบ $\sigma_1^2,\sigma_2^2$ & $Z$\\
$\mu_1-\mu_2$ อิสระกัน, ไม่ทราบแต่ $\sigma_1^2=\sigma_2^2$ & $T$ pooled, $v=n_1+n_2-2$\\
$\mu_1-\mu_2$ อิสระกัน, ไม่ทราบแต่ $\sigma_1^2\ne\sigma_2^2$ & $T$ Welch, $v$ ตามสูตร Welch--Satterthwaite\\
$\mu_D$ (ตัวอย่างสัมพันธ์กัน/วัดซ้ำ), $n<30$ & $T,\ v=n-1$\\
$\mu_D$ (ตัวอย่างสัมพันธ์กัน/วัดซ้ำ), $n\ge30$ & $Z$\\
$p$ (สัดส่วน 1 ประชากร) & $Z$ เสมอ\\
$p_1-p_2$ (ผลต่างสัดส่วน 2 ประชากร) & $Z$ เสมอ (ใช้ $\hat p_1,\hat p_2$ แยกกัน ไม่ pool)\\
\bottomrule
\end{tabular}
\end{center}

\begin{formulabox}
\textbf{ลำดับการคิดทุกครั้งที่เจอโจทย์การประมาณค่า (5 ขั้นตอน):}
\begin{enumerate}
\item \textbf{ทำความเข้าใจโจทย์} -- พารามิเตอร์ที่ถามคือค่าเฉลี่ยหรือสัดส่วน, 1 กลุ่มหรือ 2 กลุ่ม, ถ้า 2 กลุ่มเป็นอิสระหรือสัมพันธ์กัน, และโจทย์ให้ $\sigma$ มาหรือไม่
\item \textbf{กำหนดระดับความเชื่อมั่น} $(1-\alpha)$ ตามที่โจทย์ระบุ แล้วหา $\alpha/2$
\item \textbf{เลือกตัวสถิติที่ใช้} -- ใช้ตารางสรุปด้านบนตัดสินว่าเป็นกรณีใด (ทราบ/ไม่ทราบ $\sigma$, $n$ มากหรือน้อยกว่า 30)
\item \textbf{คำนวณ} -- หาค่าประมาณแบบจุดก่อน แล้วหาค่าวิกฤตจากตาราง คำนวณค่าคลาดเคลื่อนมาตรฐานและค่าคลาดเคลื่อนสูงสุดทีละขั้น สุดท้ายจึงบวก-ลบเพื่อได้ช่วง
\item \textbf{สรุปคำตอบ} -- รายงานช่วง $(L,U)$ พร้อมหน่วย และถ้าช่วงไม่คร่อมศูนย์ (สำหรับผลต่าง) ให้ระบุทิศทางของความแตกต่างประกอบด้วยเสมอ
\end{enumerate}
\end{formulabox}

\end{document}
